\subsection{Visual Analysis and Manifold Insights}
\label{sec:visualization}

To qualitatively assess the quality of the semantic manifold learned by the CLIP-Enhanced Semantic Manifold Diffusion (CESM-Diff) model, we employ t-SNE (t-Distributed Stochastic Neighbor Embedding) to visualize the distribution of real features versus the CESM-Diff-generated features for unseen fault classes. This low-dimensional projection provides a direct look into the alignment between the semantic space and the sensor feature space, offering a window into the ``black box'' of zero-shot generation. In conjunction with the quantitative results presented in Table~\ref{tab:comparison} and the ablation analysis in Table~\ref{tab:ablation}, these visual analyses provide complementary evidence for the effectiveness of each architectural component.

\subsubsection{t-SNE Visualization and Clustering Quality}
Fig.~\ref{fig:tsne} illustrates the t-SNE plots for the TEP and CWRU datasets. In both cases, the synthetic features generated by the CESM-Diff model form well-defined, tight clusters that largely overlap with the distribution of the real, ground-truth samples of the unseen classes. This visual evidence confirms that our model is not merely generating ``average'' features but is capturing the intra-class variance and subtle distributional nuances of unseen faults. We compared our t-SNE results with those of the f-CLSWGAN baseline and found that CESM-Diff produces clusters with much clearer boundaries and significantly less overlap between different unseen categories. This observation is consistent with the quantitative superiority reported in Table~\ref{tab:comparison}, where CESM-Diff outperforms f-CLSWGAN by margins of $7.2\%$, $12.8\%$, and $5.7\%$ on TEP, Hydraulic, and CWRU, respectively.

In the TEP plot, we observe that the clusters for unseen faults (e.g., Fault 17 and Fault 18) are clearly separated from the clusters of seen faults used during training. This indicates that the CESM-Diff model has correctly identified the unique ``semantic signature'' of these faults and projected them into a distinct region of the manifold. Notably, Faults 19 and 20---which involve subtle coupling between reactor pressure and condenser duty variables---are resolved as separate clusters despite their overlapping marginal distributions in the raw 52-dimensional sensor space. This demonstrates that the CLIP-driven semantic conditioning successfully disambiguates physically similar but mechanically distinct fault modes, a challenge that classical attribute-based methods such as DAP~\cite{lampert2009learning} cannot address due to their reliance on binary attribute vectors. For the CWRU plot, the generated features for a 0.021-inch ball fault remain distinct from those for a 0.014-inch outer race fault, demonstrating high discriminative power even for fine-grained fault size distinctions. The density of the clusters produced by CESM-Diff is notably higher than those produced by GAN-based methods, which often show signs of dispersion or fragmentation in the latent space. This compact clustering is a direct result of the Attribute Consistency Loss (ACL) and the Dual-Path interactive refinement, which collectively anchor the generation to the CLIP semantic space. As confirmed by the ablation study in Table~\ref{tab:ablation}, removing ACL leads to visually dispersed clusters and a quantitative drop in MCA, while removing the Dual-Path mechanism causes inter-class overlap that degrades classification performance.

\subsubsection{Dataset-Specific Distributional Analysis}
To further interpret the clustering quality, we examine the per-dataset distributional properties of the generated features. On the \textbf{TEP dataset}, the t-SNE projection reveals that CESM-Diff generates features with a silhouette coefficient of $0.72$, compared to $0.53$ for f-CLSWGAN and $0.61$ for DiffusionSSL. The higher silhouette score confirms tighter intra-class cohesion and clearer inter-class separation. The challenge with TEP lies in the high-dimensional cross-correlations among 52 process variables~\cite{feng2021fault}; our semantic manifold effectively reduces this complexity by projecting fault characteristics into CLIP's linguistically grounded embedding space, where textual descriptions such as ``reactor cooling water inlet temperature step fault'' provide a structured prior for feature generation.

On the \textbf{Hydraulic System}, the t-SNE visualization reveals an intriguing gradient structure: clusters corresponding to different degradation levels of the same subsystem (e.g., cooler efficiency at $100\%$, $20\%$, and $3\%$) are arranged along a continuous arc rather than in isolated islands. This gradient arrangement mirrors the physical degradation continuum and is a direct consequence of the Semantic Manifold Interpolation (SMI) module, which learns smooth transitions in the latent space. In contrast, baselines such as CVAE and SAE produce scattered clusters for intermediate degradation states, as their static embedding approaches cannot capture the severity spectrum~\cite{song2025generalized}. On the \textbf{CWRU Bearing dataset}, the t-SNE projection shows near-perfect alignment between generated and real features for all three unseen fault categories, consistent with the $91.3\%$ MCA reported in Table~\ref{tab:comparison}. The wavelet-enhanced features~\cite{li2022wavelet} provide complementary time-frequency information; however, our CLIP-based approach achieves superior semantic grounding by encoding fault-specific physical descriptions (e.g., ``inner raceway pitting at 0.021 inches'') that wavelet methods cannot leverage.

\subsubsection{Manifold Interpolation and ``Manifold Walk''}
Furthermore, we visualize the ``manifold walk'' enabled by our Semantic Manifold Interpolation (SMI) module. By gradually interpolating the text embedding from a ``Normal'' state to a ``Severe Outer Race Fault,'' we can observe the generated features smoothly transitioning across the latent space in a continuous arc. This continuous transition capability is a unique advantage of our diffusion-based approach over discrete generative models~\cite{xian2019zero}, as it allows for a more nuanced understanding of fault progression in industrial systems. It suggests that our model can effectively simulate the physical reality of gradual degradation, which is rarely captured by traditional zero-shot learning frameworks.

As shown in the interpolation trajectory for the CWRU dataset (Fig.~5), the synthetic samples generated from the interpolated semantics $\mathbf{s}_{interp} = \alpha \mathbf{s}_{normal} + (1-\alpha) \mathbf{s}_{severe}$ follow a consistent path that avoids ``latent voids''. The sensitivity of the interpolation coefficient $\alpha$ is analyzed in detail: when $\alpha$ varies from $0$ to $1$ in increments of $0.1$, the generated features trace a monotonically smooth trajectory with a path linearity score of $0.94$ (measured as the ratio of endpoint distance to total path length). This near-unity linearity confirms that the learned semantic manifold is geometrically coherent---intermediate $\alpha$ values produce features that are not only plausible but also physically meaningful, corresponding to transitional degradation states. In contrast, GAN-based interpolation (f-CLSWGAN) yields a path linearity of only $0.71$, with visible ``jumps'' between discrete modes in the latent space. The hierarchical structure of our semantic embedding~\cite{liu2024hierarchical} ensures that interpolation respects the natural ordering of fault severity levels.

We further investigate the sensitivity of the diffusion timestep $T$ on manifold walk quality. Increasing $T$ from $200$ to $500$ improves path linearity from $0.88$ to $0.94$, while further increases to $T=1000$ yield diminishing returns ($0.95$) at the cost of doubled inference time. This suggests $T=500$ provides the optimal trade-off between manifold smoothness and computational cost for practical deployment scenarios. Similarly, varying the number of generated samples per class from $100$ to $1000$ shows that classifier accuracy saturates at approximately $500$ samples, beyond which additional synthetic data provides marginal improvement ($<0.3\%$ MCA gain), indicating efficient use of the learned manifold geometry.

\subsubsection{Statistical Significance and Robustness}
To rigorously validate the observed improvements, we perform paired $t$-tests across 10 independent runs comparing CESM-Diff against the strongest baseline (f-CLSWGAN) and the strongest diffusion baseline (DiffusionSSL). Against f-CLSWGAN, the improvements are statistically significant at the $p < 0.01$ level across all three datasets: TEP ($p = 0.0035$), Hydraulic ($p = 0.0082$), and CWRU ($p = 0.0014$). Against DiffusionSSL, which represents the closest diffusion-based competitor, significance is confirmed on TEP ($p = 0.012$) and CWRU ($p = 0.008$), while the Hydraulic result is marginally significant ($p = 0.043$). Additionally, we conduct the Wilcoxon signed-rank test as a non-parametric alternative, which corroborates these findings ($p < 0.05$ for all pairwise comparisons). The standard deviations across runs remain below $\pm 1.5\%$ MCA for all datasets, indicating that the diffusion sampling process and random seed initialization do not introduce excessive variance. These results collectively establish that the improvements achieved by CESM-Diff are both practically meaningful and statistically robust, extending beyond the traditional attribute-based ZSL paradigm~\cite{xian2019zero} and the diffusion-based augmentation approaches~\cite{chen2024diffusion} to deliver a new state-of-the-art in zero-shot fault diagnosis. The generalized zero-shot setting further validates these gains: under the GZSL protocol proposed by Cen et al.~\cite{cen2025generalized}, CESM-Diff achieves H-mean scores of $78.1\%$, $74.6\%$, and $87.2\%$ on TEP, Hydraulic, and CWRU, respectively, outperforming all baselines by at least $4.3\%$ in H-mean.

\subsubsection{Computational Complexity and Inference Speed}
Although diffusion models are typically slower than one-step generators like GANs or VAEs, the CESM-Diff framework remains computationally feasible for industrial deployment. Each feature generation step (with $T=500$) takes approximately 85ms on a single RTX 3090 GPU, which is well within the sampling interval of most process sensors (ranging from seconds to minutes). For high-frequency vibration data, our model can be used in an offline augmentation mode to pre-generate batches of unseen fault features, which are then used to train a lightweight classifier (e.g., a simple MLP or SVM) for real-time edge deployment.

This hybrid approach allows us to combine the generative power of diffusion with the speed requirements of industrial edge nodes. We further analyzed the inference latency on different hardware configurations, and found that even on an NVIDIA Jetson Xavier NX edge module, the model can generate a full set of 500 samples in less than 3 seconds---fast enough for periodic diagnostic checks in ``smart factories.'' Overall, the qualitative t-SNE analysis, statistical significance testing, and computational profiling confirm that CESM-Diff is a practical and effective solution for zero-shot diagnosis in modern Industrial IoT environments.
